% !TeX root = ../../Book3.tex
% 纲要标题：### G.4 正性证书与 Positivstellensatz

\section{正性证书与 Positivstellensatz}
\label{b3:sec:G04}

约束集合上的正性可以通过约束多项式参与的平方和恒等式来验证。是否允许约束间的乘积、是否具有紧性或 Archimedes 条件，会决定证书的形式与存在范围。

% 纲要细目（与计划纲要同步）：
% * 预序、二次模及多项式不等式系统；以恒等式作为不可行性或正性的证书
% * Krivine–Stengle 型 Positivstellensatz 的一个明确版本；等式约束、严格不等式和非严格不等式分别处理
% * Schmüdgen／Putinar 型结果只作比较：紧性、Archimedean 条件及严格正性不能无条件互换
%
% #### G.4.1 预序、二次模与多项式不等式证书
% #### G.4.2 Krivine–Stengle 型 Positivstellensatz
% #### G.4.3 Schmüdgen 与 Putinar 型结果的假设比较


% 正文按纲要完成；所列定理给出证明或证明概要。

\subsection{预序、二次模与多项式不等式证书}
\label{b3:subsec:G04:01}

设约束集合为 \(K=\{x\in\mathbb R^n:g_1(x)\ge0,\ldots,g_s(x)\ge0\}\)。若能写出 \(f=\sigma_0+\sum_i\sigma_i g_i\)，各 \(\sigma_i\) 都为平方和，就已经证明 \(f\) 在 \(K\) 上非负。再允许约束之间的乘积，会得到更广的一类证书。

\begin{definition}\label{b3:def:preordering-quadratic-module}
设 \(R\) 为实闭域，\(A=R[x_1,\ldots,x_n]\)，记 \(\Sigma A^2\) 为有限平方和的集合。给定 \(g_1,\ldots,g_s\in A\)，集合
\[
M(g)=\Sigma A^2+\sum_{i=1}^s\Sigma A^2g_i,
\qquad
T(g)=\sum_{\varepsilon\in\{0,1\}^s}\Sigma A^2g_1^{\varepsilon_1}\cdots g_s^{\varepsilon_s}
\]
分别称为它们生成的\term{二次模}[quadratic module]与\term{预序}[preordering]。
\end{definition}

这里的二次模是含 \(1\)、对加法与乘平方封闭的子集，不是通常意义的 \(A\) 模。预序还对乘法封闭：把乘积中 \(g_i\) 的偶次幂吸收到平方和系数即可。二者中每个元素都在 \(K\) 上非负。比如在 \([-1,1]\) 上，\(1-x=(1-x)^2/2+(1-x^2)/2\) 是由约束 \(1-x^2\ge0\) 给出的精确证书。

\subsection{Krivine–Stengle 型 Positivstellensatz}
\label{b3:subsec:G04:02}

\begin{theorem}[Krivine--Stengle 型不可行性证书]\label{b3:thm:real-infeasibility-certificate}
设 \(R\) 为实闭域，\(f_1,\ldots,f_r,g_1,\ldots,g_s\in R[x_1,\ldots,x_n]\)，并令 \(I=(f_1,\ldots,f_r)\)、\(T(g)\) 为 \(g_1,\ldots,g_s\) 生成的预序。则
\[
\{x\in R^n:f_i(x)=0\ (1\le i\le r),\ g_j(x)\ge0\ (1\le j\le s)\}=\varnothing
\quad\Longleftrightarrow\quad -1\in I+T(g).
\]
\end{theorem}

\begin{proof}
若 \(-1\in I+T(g)\)，在任何假定的解上求值，理想部分为零、预序部分非负，不可能得到 \(-1\)。证明反向的逆否命题。设 \(-1\notin I+T(g)\)，在 \(B=R[x]/I\) 中取由 \(g_j\) 的像生成的预序，利用 Zorn 引理把它扩为不含 \(-1\) 的极大预序 \(P\)。预序包含全部平方、对加法与乘法封闭；链的并仍有这些性质且不含 \(-1\)。

说明极大性使它成为一个带素支集的正锥。若 \(a,-a\) 都不在 \(P\)，加入它们中的任一个所得预序分别为 \(P+aP\)、\(P-aP\)，因为 \(a^2\in P\)。两者都不能再扩大 \(P\)，于是有
\(-1=u+av\)、\(-1=w-az\)，其中 \(u,v,w,z\in P\)。相乘整理得
\(1+u+w+uw+a^2vz=0\)，即 \(-1\in P\)，矛盾。因此 \(P\cup(-P)=B\)。集合 \(\mathfrak q=P\cap(-P)\) 是理想：乘任意元时，可根据该元或其负元属于 \(P\) 来验证封闭性。

它还是素理想。若 \(ab\in\mathfrak q\) 而 \(a,b\notin\mathfrak q\)，更换符号后可令 \(a,b\in P\)、\(-a,-b\notin P\)。极大性给出 \(-1=u-av\)、\(-1=w-bz\)，故
\(abvz=(1+u)(1+w)\)。左边属于 \(\mathfrak q\)，把它的负元与 \(u+w+uw\in P\) 相加，再次得到 \(-1\in P\)，矛盾。

于是 \(B/\mathfrak q\) 是有序整环，其分式域 \(L\) 带有延拓的序，且各 \(g_j\) 的像非负。\(R\) 在 \(L\) 中仍保持原序，因为实闭域的正元都是非零平方。取 \(L\) 的一个实闭包 \(L^{\mathrm{rc}}\)，\(x_i\) 的像在这个实闭扩域内满足原系统。\cref{b3:thm:real-projection}给出的量词消去说明带 \(R\) 系数的存在公式在 \(R\subseteq L^{\mathrm{rc}}\) 中真值相同：消去后只剩 \(R\) 中元素的符号，而嵌入保持这些符号。因此原系统在 \(R\) 中也有解。
\end{proof}

实闭包的构造见第一卷附录 E。所证版本在 \(R=\mathbb R\) 时是\cite[Section 7.4, Theorem 7.5]{Sturmfels2002PolynomialEquations}没有严格不等式的情形；证明没有调用实零点定理，所以稍前使用此结果推导实根理想公式不会产生循环。

\ProseParagraphBreak
严格不等式和不等条件也可以转换后使用该定理。例如 \(h(x)\ne0\) 等价于存在 \(z\) 使 \(zh(x)-1=0\)；\(h(x)>0\) 等价于存在 \(z\) 使 \(h(x)z^2-1=0\)。后一等价使用正实数有平方根。转换增加变量，却保留了精确可行性。

\begin{example}\label{b3:exm:inconsistent-inequalities}
条件 \(x\ge0\)、\(-x-1\ge0\) 不可能同时成立，恒等式 \(-1=x+(-x-1)\) 就是预序证书。若另有等式 \(y-x^2=0\) 并要求 \(y<0\)，直接利用 \(y=x^2\) 已知矛盾；上述变量替换还给出等式证书：加入 \((-y)z^2-1=0\) 后，
\[
-1=(-yz^2-1)+z^2(y-x^2)+(xz)^2.
\]
最后一项为平方，前两项属于等式理想。
\end{example}

\subsection{Schmüdgen 与 Putinar 型结果的假设比较}
\label{b3:subsec:G04:03}

\begin{theorem}[紧集上的严格正性]\label{b3:thm:compact-positivity-certificates}
设 \(g_1,\ldots,g_s,f\in\mathbb R[x_1,\ldots,x_n]\)，\(K=\{x\in\mathbb R^n:g_i(x)\ge0\text{ 对每个 }i\}\)，并记 \(T(g)\)、\(M(g)\) 分别为 \(g_1,\ldots,g_s\) 生成的预序与二次模。
\begin{enumerate}
\item 若 \(K\) 紧且 \(f>0\) 于 \(K\)，则 \(f\in T(g)\)；这是 Schmüdgen 型结论。
\item 若存在 \(N>0\) 使 \(N-\sum_i x_i^2\in M(g)\)，且 \(f>0\) 于 \(K\)，则 \(f\in M(g)\)；这是 Putinar 型结论。
\end{enumerate}
\end{theorem}

\begin{proofsketch}
先证第二项，采用\cite[Section 2, Theorem 6, Lemma 8 and proof of Theorem 3]{Schweighofer2005Optimization}的构造。以下置 \(A=\mathbb R[x_1,\ldots,x_n]\)。记 \(b=N-\sum_i x_i^2\in M(g)\)，\(a=N+1/4\)，并令
\[
\Delta=\{y\in\mathbb R_{\ge0}^{2n}:\textstyle\sum_jy_j=2na\},
\qquad \ell(y)_i=(y_i-y_{n+i})/2,
\qquad C=\ell(\Delta).
\]
这是紧的单纯形及其线性像。把各 \(g_i\) 乘以适当正常数，使它们在 \(C\) 上都不超过一；这不改变 \(K\) 或 \(M(g)\)。可以选择 \(c\ge0\) 及充分大的整数 \(k\)，使
\[
q=f-c\sum_i(1-g_i)^{2k}g_i>0\quad\text{于 }C.
\]
理由是：在紧集 \(\{x\in C:f(x)\le\varepsilon\}\) 上，取 \(\varepsilon>0\) 使它与 \(K\) 不交，则至少有一个 \(g_i\le-\delta\)，其中 \(\delta>0\) 可统一选择；这个负项的扣除可使 \(f\) 增大。对于 \(0\le t\le1\)，有 \(t(1-t)^{2k}\le1/(2k+1)\)，所以所有非负项造成的损失随 \(k\) 增大而统一趋于零。先选 \(c\)，再选 \(k\)，便使上式在这个坏集及其补集都为正；坏集为空时取 \(c=0\)。

将 \(q\) 经 \(\ell\) 提升到 \(\Delta\)，再用 \(\sum y_j/(2na)\) 把各项补齐次数，得到在非负象限除原点以外严格正的齐次式 \(Q(y)\)。P\'olya 的系数引理保证，对充分大的 \(e\)，
\(Q(y)(\sum y_j/(2na))^e\) 的系数全非负。引理的证明是比较其按多项式系数归一化后的系数与 \(Q(\alpha/|\alpha|)\)：前者中的下降阶乘与普通幂之差为一致的 \(O(1/e)\)，而 \(Q\) 在标准单纯形上的最小值严格大于零。本概要省略这个有限系数估计的展开，详见上述 Theorem 6。

代入 \(y_i=a+x_i\)、\(y_{n+i}=a-x_i\) 后，\(\sum y_j/(2na)=1\)，所得式正是 \(q\)。而
\[
a\pm x_i=b+\sum_{j\ne i}x_j^2+(x_i\pm1/2)^2
\in T(b)=\Sigma A^2+b\Sigma A^2\subseteq M(g).
\]
\(T(b)\) 对乘法封闭，故上述非负系数表示给出 \(q\in T(b)\)。此前减去的每项属于 \(M(g)\)，所以 \(f\in M(g)\)。这里没有把一般二次模当作对乘法封闭的预序。

第一项还须从紧性得到 \(T(g)\) 内的球约束。以下把这个步骤写出。记 \(\Sigma=\sum_i x_i^2\)，选 \(s>0\) 使 \(h=s-\Sigma>0\) 于 \(K\)。系统 \(g_i\ge0,-h\ge0\) 不可行；\cref{b3:thm:real-infeasibility-certificate}给出 \(h t_1=1+t_0\)，其中 \(t_0,t_1\in T(g)\)。由恒等式
\[
h\bigl(1+(h-2)^2t_1\bigr)
=1+(h-3/2)^2+3/4+(h-2)^2t_0
\]
得到 \(h(1+v)=1+u\)，其中 \(u,v\in T(g)\)。令 \(T(h)=\Sigma A^2+h\Sigma A^2\)。对任意多项式 \(p\)，均存在常数 \(t\ge0\) 使 \(t\pm p\in T(h)\)：具有此性质的多项式对加法、乘法封闭，乘法用
\[
tu\pm pq=\tfrac12\bigl((t+p)(u\pm q)+(t-p)(u\mp q)\bigr)
\]
验证；而 \(s+1/4\pm x_i=h+\sum_{j\ne i}x_j^2+(x_i\pm1/2)^2\) 说明它包含各变量。

因此可取 \(t\ge0\) 使 \(t-sv\in T(h)\)。由于 \(1+v\in T(g)\) 且 \(h(1+v)\in T(g)\)，有
\((1+v)(t-sv)\in T(g)\)。另外 \(h+sv=h(1+v)+v\Sigma\in T(g)\)。将这两式与平方 \(s(t/(2s)-v)^2\) 相加，得到
\[
\bigl(s+t+t^2/(4s)\bigr)-\Sigma\in T(g).
\]
于是把已证第二项用于二次模 \(T(g)\)，即得 \(f\in T(g)\)。此球约束构造也说明为什么这一步依赖预序的乘法封闭性。所用弱紧性结论及两项的关系可对照\cite[Section 1, Theorems 1 and 3]{Schweighofer2005Optimization}。
\end{proofsketch}

第二个条件称为二次模的 Archimedean 条件的一种等价形式，它显式给出球约束，因而蕴含 \(K\) 紧；仅知 \(K\) 紧不能直接删去这个条件。两个结论的证书格式也不同：预序允许 \(g_ig_j\)，二次模只允许单个生成约束。

\ProseParagraphBreak
严格正性不能任意换成非负性。例如令 \(g=(1-x^2)^3\)，则 \(K=[-1,1]\)，\(M(g)\) 满足上述条件：由定理第一条及只有一个生成元时 \(T(g)=M(g)\)，可知 \(2-x^2\in M(g)\)。然而 \(1-x^2\notin M(g)\)。若有 \(1-x^2=\sigma_0+\sigma_1(1-x^2)^3\)，在 \(x=1\) 得 \(\sigma_0(1)=0\)，于是组成 \(\sigma_0\) 的每个平方的底数在 \(1\) 为零，故 \(\sigma_0'(1)=0\)。右边导数在 \(1\) 为零，左边却为 \(-2\)。这个简单的消失阶数检验解释了为何定理要保留 \(f>0\)。
